% Source: sources/returns/S6_04_MAPS_MANIFOLDS/sections/manifolds.tex:1671–1799
\paragraph{Theorem D.2.9 (Proved: all orders $4k$ in the Liu--Rollenske configuration).}
Use the $III^*,II,I_1$ surface and the unchanged type-$II$
filling of Liu--Rollenske, with the degree-one linearization zero
at $I_1$ and the primitive hexagonal filling there.
In their homology basis $(e_1,e_2,e_3,e_4)$ put
\[
 \delta=-e_1+2e_3,\qquad \ell=e_4,\qquad
 m=4k,\qquad v=n\ell+r\delta,
 \quad k>0,\quad\gcd(n,4k)=1.
\]
The logarithmic transformation at $III^*$ exists for every
such integer datum, and its total space satisfies
\[
 \pi_1(X)=\mathbb Z/|n|,\qquad
 H_0=H_6=\mathbb Z,\quad H_5=0,\quad
 H_1=H_2=H_3=H_4=\mathbb Z/|n|.
\]
Thus $n=\pm1$ gives $S^6$ for every $k>0$. In particular,
the Liu--Rollenske configuration also answers Viaclovsky's
Question 4 affirmatively. The type-$II$ fibre is unchanged
in this statement.
\begin{proof}
The source's good-reduction construction and section-preserving
comparison at $III^*$ give the analytic model. The degree-$k$
pullback and free affine quotient argument of D.2.6 applies
with invariant covector $e_4^*$, which takes value $n$ on $v$.
It proves existence, smoothness, properness and the relation
$g^m=v$ without restricting $k$.

The extra local input, absent from D.2.6, is the unchanged
type-$II$ filling. By \cite[Lemma 3.10, Proposition 3.11]{manLR26},
its marked neighbourhood is homeomorphic to the type-$II$
elliptic neighbourhood times $T^2$, with the section preserved.
The cuspidal rational curve is homeomorphic to $S^2$.
The section evaluates as one on both this central two-cycle
and the nearby elliptic curve, so the degree-two specialization
is primitive. The product description makes its two normal
tori precisely the meridian swept by each of the two transverse
$T^2$ loops: each bounds the section disc times that loop.
K\"unneth, or the relative sequence of this product pair,
shows that these generate the entire rank-two kernel of
the boundary map in degree two. This justifies the integral
attachment used below without treating the cusp as a normal
crossing curve.

Here are all matrices and primitive certificates. The source
uses $T_\infty T_1T_0=I$ in homology. We order its matrices
as $(A_0,A_1,A_2)=(T_0,T_\infty,T_1)$, so their product is $I$:
\[
 \begin{gathered}
 A_0=\begin{pmatrix}-1&-1&-1&0\\2&1&1&0\\0&0&1&0\\0&0&0&1\end{pmatrix},\qquad
 A_1=\begin{pmatrix}1&0&0&-1\\-2&0&-1&1\\2&1&2&1\\0&0&0&1\end{pmatrix},\\
 A_2=\begin{pmatrix}1&1&0&1\\-2&0&-1&-1\\0&-1&1&-1\\0&0&0&1\end{pmatrix}.
 \end{gathered}
\]
These are \cite[Proposition 3.13]{manLR26} in its actual
integral marking. The cusp's primitive vanishing lattice is
$\langle a,\delta\rangle$, where $a=e_1-e_2-e_3$.
Its monodromy-normal closure contains
$A_2a=-e_2$, and hence all of
$\langle e_1,e_2,e_3\rangle$. The two unchanged meridians
are killed by the extending section. The sphere relation kills
the third, and $g^m=v$ then leaves only $ne_4=0$.
Surjectivity of the boundary maps on fundamental groups gives
$\pi_1(X)=\mathbb Z/|n|$.

Set $D=[A_0-I\ \ A_1-I]$ and, in $\Lambda^2$, define
\[
 j_0(x)=(x,0),\quad j_1(x)=(0,x),\quad
 j_2(x)=-(A_1x,x).
\]
The cusp fixed basis is $(a,\delta)$ and the type-$II$ fixed
basis is $(\delta,e_4-e_2)$. Form
\[
 K=[j_1(a)\ j_1(\delta)\ j_2(\delta)\
       j_2(e_4-e_2)\ j_0(\ell)].
\]
Then $DK=0$, $\operatorname{rank}D=3$, and the $5\times5$
minor on rows $(1,2,4,5,6)$ is $1$.
Moreover $j_0(\delta)=-K_2-K_3$.
The exterior matrix $[\bigwedge^2A_i-I]_{i=0}^2$ has rank
five and the minor on rows $(1,2,3,4,5)$ and columns
$(2,4,5,6,11)$ is $1$. Its primitive cokernel is detected by
\[
 \omega=2e_1^*e_2^*-e_2^*e_3^*-e_3^*e_4^*.
\]
Thus $H_2(V)$ is free of rank six, split by a fibre class
$s=-e_3e_4$, of $\omega$-value one, and the five zero-value
sweeps $K$. The four normal tori of the two unchanged fillings
are exactly its first four sweep columns, by the preceding
product argument and the rank-two argument in D.2.6.
After quotienting by them the group is
$Q=\mathbb Z\langle s,t\rangle$, $t=j_0(\ell)$.

The finite linear coinvariants satisfy $e_1=e_2=0$, so
$\delta=2e_3$. The finite reduced fibre therefore has free
$H_1$ with basis $e_3,z$, where
\[
 mg=ne_4+2r e_3,\qquad z=\alpha e_4+\beta g,
 \qquad m\alpha+n\beta=1.
\]
Direct calculation gives
\[
 \mathcal N_me_3=2k\delta,\quad
 \mathcal N_me_4=m\ell,\quad
 \omega(e_3,v)=-n,\quad \omega(e_4,v)=2r.
\]
The $\delta$ sweep dies in $Q$, and the normal-circle formula
therefore reduces the last map to
\[
 \begin{pmatrix}n&-2r\alpha\\0&1\end{pmatrix}.
\]
All boundary $H_1$ groups are free of total rank nine, and their
target is free of rank nine with finite cokernel
$\mathbb Z/|n|$; the degree-one kernel is zero.
Every central fibre is irreducible and a multiple is principal,
so the degree-two boundary maps are onto as in D.2.5.
It follows that $H_2(X)=\mathbb Z/|n|$.
The unchanged type-$II$ piece and the finite quotient have
Euler number zero, whereas the hexagonal fibre has Euler
number two. Duality yields all the other displayed groups.
Sphere recognition when $|n|=1$ is as in D.2.6.

The independent script
\texttt{code/manifolds\_liu\_rollenske.py} checks the two
unit minors, the norm identities, and 312 full $6\times6$
Smith forms for the same $(k,n,r)$ bounds as D.2.6.
The proof establishes the formula without those bounds.
\end{proof}
